\section{Feeder Co-Simulation Methodology}
% provenance: mis_01KXT2DXXY862RPX7YATTRBDWB (M4 spec); dec_01KXT27V1Y4ZR6ZRAST3H94WC7 (testbed decision).
% provenance: dec_01KXT4C6ZPYD1Z92DH6CN1KKM3 (BR_phys unit = ANSI C84.1 excursion + trip flag).

The physical consequence radius is measured by a feeder co-simulation that translates an
authorized-control reach from the analysis engine into a malicious dispatch program and runs
it on a public distribution feeder. We take the IEEE 8500-node feeder as the primary
benchmark and the PNNL 9500-node system as an advanced validation tier. Time-series inputs
come from public solar and load datasets. This part of the work is methodologically routine
relative to the false-data-injection literature on distribution systems, so we build it to be
unimpeachable rather than novel.

Two choices harden it against the standard objections to simulation studies. First, we cross-
check OpenDSS against an independent GridLAB-D solver on the base case and on at least one
attack case, following the precedent of the PNNL 9500 report, and we report any divergence
rather than suppressing it. Second, credential state gates dispatch inside the co-simulation
through a HELICS federation, so that a denied or narrowed credential actually prevents a
command in simulation rather than being applied after the fact. The attack families are
volt-var curve distortion, volt-watt curve distortion, export-limit distortion, and
curtailment spam. The metrics are the count of nodes violating the ANSI~C84.1 bands, maximum
voltage deviation, the area under the violation curve over time, regulator tap operations,
unnecessary curtailment, energy not served, and recovery time to steady limits, with the
per-unit voltage excursion as the primary continuous measure and a protection-trip flag as
the discrete outcome.

We do not tune feeder parameters to amplify impact; we report what the standard models
produce. The clean base case, the four attack families, the eight metrics, and the
cross-solver agreement are the acceptance criteria for this component. As with the credential
service, a first benchmark run on the IEEE~8500-node feeder appears in the preliminary evaluation;
the full sweep on both feeders, with the GridLAB-D cross-check and seasonal operating states,
follows the pre-registered protocol chosen in advance.
